\documentclass{article}
\usepackage{amsmath,amssymb,amsthm}
\usepackage{mathtools}
\usepackage{hyperref}

\newtheorem{theorem}{Theorem}[section]
\newtheorem{proposition}[theorem]{Proposition}
\newtheorem{lemma}[theorem]{Lemma}
\theoremstyle{remark}
\newtheorem{remark}[theorem]{Remark}

\DeclareMathOperator{\Var}{Var}
\DeclareMathOperator{\dist}{dist}

\title{A counterexample to uniform $L^2$ stability\\
       in the Brascamp--Lieb variance inequality}
\author{}
\date{}

\begin{document}
\maketitle

\begin{abstract}
We show that the sharp uniform $L^1$ stability estimate for the Brascamp--Lieb
variance inequality cannot be strengthened to an $L^2$ estimate with a constant
depending only on the dimension. Already in dimension one we exhibit, for all
sufficiently large $S$, a smooth strictly convex potential $\phi_S$ and a
locally Lipschitz test function $f_S$ such that
\[
  \frac{\dist_{L^2(\rho_{\phi_S})}(f_S,\mathcal O_{\mathrm{BL},\phi_S})^2}
       {\delta_{\mathrm{BL},\phi_S}(f_S)}
  = S\bigl(1+O(S^{-1})\bigr)\to+\infty.
\]
Equivalently, the constant in any putative inequality
$\dist_{L^2}\le C\,\delta_{\mathrm{BL}}^{1/2}$ must blow up at rate at least
$\sqrt S$.
\end{abstract}

\section{Statement}

Let $\phi:\mathbb R^d\to\mathbb R\cup\{+\infty\}$ be an essentially continuous
convex function with $0<\int e^{-\phi}<\infty$. Write
\[
  Z_\phi := \int_{\mathbb R^d} e^{-\phi}\,dx, \qquad
  d\rho_\phi := Z_\phi^{-1}e^{-\phi}\,dx.
\]
For smooth strictly convex $\phi$, the \emph{Brascamp--Lieb deficit} is
\[
  \delta_{\mathrm{BL},\phi}(f)
  =
  \int_{\mathbb R^d}\bigl\langle (D^2\phi)^{-1}\nabla f,\nabla f\bigr\rangle\,d\rho_\phi
  -\Var_{\rho_\phi}(f),
\]
defined for any locally Lipschitz $f\in L^2(\rho_\phi)$. The
Brascamp--Lieb inequality \cite{BL76} asserts $\delta_{\mathrm{BL},\phi}(f)\ge 0$,
with equality cases
\[
  \mathcal O_{\mathrm{BL},\phi}
  =
  \{a\cdot \nabla\phi+b:\ a\in\mathbb R^d,\ b\in\mathbb R\}.
\]
When using $L^2(\rho_\phi)$ distances, this affine family is understood through
its members that belong to $L^2(\rho_\phi)$; for the explicit potentials
constructed below all the listed generators are in $L^2$.

\begin{theorem}\label{thm:l2-fails}
There is no finite constant $C$ such that
\begin{equation}\label{eq:would-be-stability}
  \dist_{L^2(\rho_\phi)}(f,\mathcal O_{\mathrm{BL},\phi})
  \le C\,\delta_{\mathrm{BL},\phi}(f)^{1/2}
\end{equation}
holds simultaneously for every smooth strictly convex
$\phi:\mathbb R\to\mathbb R$ with $\int e^{-\phi}<\infty$ and every locally
Lipschitz $f\in L^2(\rho_\phi)$. Consequently, no dimension-only constant
$C=C(d)$ can exist in any dimension $d\ge 1$.
\end{theorem}

\subsection*{Spectral interpretation}

In dimension one, the Brascamp--Lieb Dirichlet form is
\begin{equation}\label{eq:BL-form}
  \mathcal E_\phi(f,g)
  := \int_{\mathbb R}\frac{f'g'}{\phi''}\,d\rho_\phi,
  \qquad
  \mathcal E_\phi(f) := \mathcal E_\phi(f,f).
\end{equation}
Here the form domain means the locally absolutely continuous functions
$f\in L^2(\rho_\phi)$ for which $\mathcal E_\phi(f)<\infty$.
After closing the form, its restriction to the orthogonal complement of
constants in $L^2(\rho_\phi)$ defines a non-negative self-adjoint operator
$L_\phi$ via
$\mathcal E_\phi(f,g)=\langle L_\phi f,g\rangle_{L^2(\rho_\phi)}$, and the
Brascamp--Lieb inequality is the assertion that
$\inf\,\mathrm{spec}\bigl(L_\phi\restriction\{1\}^\perp\bigr)\ge 1$.
Theorem~\ref{thm:l2-fails} is the assertion that the spectral gap
$\inf\,\mathrm{spec}\bigl(L_\phi\restriction \mathcal O_{\mathrm{BL},\phi}^\perp\bigr)-1$
\emph{cannot} be bounded below by a positive universal constant: the
construction below produces potentials for which the Rayleigh quotient on
$\mathcal O_{\mathrm{BL},\phi}^\perp$ approaches~$1$ from above. If one also
knows compactness of the resolvent for the corresponding one-dimensional
operator, this can be phrased as collapse of the next eigenvalue.

The basic optimizer direction at the lowest level is $\phi'$ itself.

\begin{lemma}\label{lem:phi-prime-eigen}
Let $\phi:\mathbb R\to\mathbb R$ be smooth and strictly convex, with
$0<\int e^{-\phi}<\infty$. Assume in addition that $\phi''$ is bounded and
$\phi'\in L^2(\rho_\phi)$; in particular, this holds for the potentials
$\phi_S$ constructed below. Then $\int\phi'\,d\rho_\phi=0$, $\phi'$ belongs to
the form domain, and
\[
  \mathcal E_\phi(\phi',g)=\langle\phi',g\rangle_{L^2(\rho_\phi)}
  \qquad\text{for every locally Lipschitz form-domain }g.
\]
In particular $\mathcal E_\phi(\phi')=\Var_{\rho_\phi}(\phi')$, and
$\phi'$ satisfies the weak eigenvalue identity at level~$1$ against locally
Lipschitz form-domain test functions.
\end{lemma}

\begin{proof}
Integrability of $e^{-\phi}$ on $\mathbb R$ together with strict convexity
forces $\phi(x)\to+\infty$ as $|x|\to\infty$, hence $e^{-\phi(\pm R)}\to 0$
as $R\to\infty$. Moreover, by convexity $\phi'$ is monotone non-decreasing,
so $\phi'(x)$ has a (possibly infinite) limit as $x\to\pm\infty$, and there
exists $x_0>0$ with $\phi'(x)\ge 0$ for $x\ge x_0$ and $\phi'(x)\le 0$ for
$x\le -x_0$ (otherwise $\phi$ could not satisfy $\phi(\pm\infty)=+\infty$).
Thus
\[
  \int_{\mathbb R}|\phi'(x)|e^{-\phi(x)}\,dx
  = -\int_{-\infty}^{-x_0}(e^{-\phi})'\,dx
    +\int_{-x_0}^{x_0}|\phi'|e^{-\phi}\,dx
    +\int_{x_0}^{\infty}-(e^{-\phi})'\,dx <\infty,
\]
and integrating $\phi'e^{-\phi}=-(e^{-\phi})'$ over the two tails and then
letting the endpoints tend to infinity gives
$\int\phi'\,d\rho_\phi=0$.

\emph{Integration by parts for general $h$.} Let
$h\in L^2(\rho_\phi)$ be locally Lipschitz and in the form domain. Since
$\phi''$ is bounded, $h'\in L^2(\rho_\phi)$. By Cauchy--Schwarz,
$h\phi' e^{-\phi}, h e^{-\phi}, h' e^{-\phi}\in L^1(\mathbb R)$.
Choose sequences $R_k\to+\infty$ and $R_k'\to-\infty$ along which
$|h(R_k)|e^{-\phi(R_k)}+|h(R_k')|e^{-\phi(R_k')}\to0$. Integration by parts on
$[R_k',R_k]$ gives
\[
  \int_{R_k'}^{R_k}h\phi' e^{-\phi}\,dx
  = -\int_{R_k'}^{R_k}h(e^{-\phi})'\,dx
  = -\bigl[h\,e^{-\phi}\bigr]_{R_k'}^{R_k}
    +\int_{R_k'}^{R_k}h' e^{-\phi}\,dx.
\]
The boundary term vanishes along the chosen sequence; passing to the limit
yields
\begin{equation}\label{eq:ipb}
  \int_{\mathbb R} h\phi' e^{-\phi}\,dx
  = \int_{\mathbb R} h' e^{-\phi}\,dx.
\end{equation}

For any locally Lipschitz form-domain $g$, applying \eqref{eq:ipb} to $h=g$
yields
\[
  \mathcal E_\phi(\phi',g)
  =\int\frac{(\phi')'g'}{\phi''}\,d\rho_\phi
  =\int g'\,d\rho_\phi
  =\int g\,\phi'\,d\rho_\phi
  =\langle\phi',g\rangle_{L^2(\rho_\phi)}.
\]
Specializing to $g=\phi'$ gives
$\mathcal E_\phi(\phi')=\|\phi'\|_{L^2(\rho_\phi)}^2
=\Var_{\rho_\phi}(\phi')$, since $\phi'$ has $\rho_\phi$-mean zero.

For the later potentials $\phi_S$, the assumptions are immediate from the
explicit formulas: $\phi''_S$ is bounded for each fixed $S$, while on the
tails $\phi'_S(x)=\eta x\pm S$ and $e^{-\phi_S(x)}$ has at least Gaussian
decay.
\end{proof}

\section{The construction}

Fix a non-negative even $C^\infty$ bump
$\kappa:\mathbb R\to[0,\infty)$ supported in $[-1,1]$ and normalized by
$\int_{\mathbb R}\kappa(u)\,du=1$. For parameters $S\ge 1$,
$\varepsilon\in(0,\tfrac14)$ and $\eta>0$ define
\begin{equation}\label{eq:phi-double-prime}
  \phi''_{S,\varepsilon,\eta}(x)
  :=\eta
  +\frac{S}{\varepsilon}\kappa\!\left(\frac{x-1}{\varepsilon}\right)
  +\frac{S}{\varepsilon}\kappa\!\left(\frac{x+1}{\varepsilon}\right),
\end{equation}
together with the normalization $\phi_{S,\varepsilon,\eta}(0)
=\phi_{S,\varepsilon,\eta}'(0)=0$. The right-hand side of
\eqref{eq:phi-double-prime} is $C^\infty$, even, and bounded below by~$\eta>0$,
so $\phi_{S,\varepsilon,\eta}$ is $C^\infty$, even, strictly convex, with
$\phi_{S,\varepsilon,\eta}(x)\ge\eta x^2/2$, hence
$\int e^{-\phi_{S,\varepsilon,\eta}}<\infty$. By construction,
$\phi'_{S,\varepsilon,\eta}$ is odd.

We require the parameters to satisfy
\begin{equation}\label{eq:scaling}
  S\varepsilon\to 0,\qquad \eta\,S^{-2}\to 0,\qquad \eta\to 0,
  \qquad\text{as }S\to\infty.
\end{equation}
For concreteness we work throughout with the explicit choice
\begin{equation}\label{eq:specific-choice}
  \varepsilon = S^{-3},\qquad \eta = S^{-4},
\end{equation}
which obeys all conditions in \eqref{eq:scaling}, and abbreviate
\[
  \phi_S:=\phi_{S,S^{-3},S^{-4}},\qquad
  \rho_S:=\rho_{\phi_S},\qquad
  Z_S:=\int_{\mathbb R}e^{-\phi_S}\,dx.
\]
All asymptotic notation $O(\,\cdot\,)$, $o(1)$ below is understood as
$S\to\infty$.

\subsection*{Notation for the three regions}

Decompose $\mathbb R$ into the central plateau, the two transition layers, and
the exterior:
\[
  C_S:=[-1+\varepsilon,1-\varepsilon],\quad
  T_S:=I_S^+\cup I_S^-,\quad
  E_S:=\mathbb R\setminus(C_S\cup T_S),
\]
where
\[
  I_S^+:=(1-\varepsilon,1+\varepsilon),\qquad I_S^-:=(-1-\varepsilon,-1+\varepsilon).
\]

\subsection*{Profile of $\phi_S$, $\phi'_S$ and the partition function}

Direct integration of \eqref{eq:phi-double-prime}, together with
$\int\kappa=1$ and the support assumption on $\kappa$, gives:
\begin{itemize}
\item On $C_S$: $\phi''_S\equiv\eta$, so $\phi'_S(x)=\eta x$ and
  $\phi_S(x)=\eta x^2/2\in[0,\eta/2]$.
\item On $I_S^+$: $\phi'_S$ is non-decreasing and
  $\int_{1-\varepsilon}^{1+\varepsilon}\phi''_S\,dx=2\eta\varepsilon+S$, so
  $\phi'_S$ ranges from $\eta(1-\varepsilon)$ at the left endpoint to
  $\eta(1+\varepsilon)+S$ at the right endpoint.
\item On $(1+\varepsilon,\infty)$: $\phi''_S=\eta$, so
  $\phi'_S(x)=\eta x+S$.
\item Symmetric statements hold on $I_S^-$ and $(-\infty,-1-\varepsilon)$.
\end{itemize}

\begin{lemma}\label{lem:phi-bounds}
Under the choice \eqref{eq:specific-choice}, for all sufficiently large $S$:
\begin{enumerate}
\item[(a)] $0\le\phi_S(x)\le\eta/2=\tfrac{1}{2}S^{-4}$ for $x\in C_S$.
\item[(b)] $0\le\phi_S(x)\le \tfrac{\eta}{2}+2\varepsilon\bigl(\eta+S\bigr)
  =O(S\varepsilon)=O(S^{-2})$ for $x\in T_S$. In particular
  $e^{\phi_S(x)}=1+O(S^{-2})$ uniformly on $T_S$.
\item[(c)] For $x>1+\varepsilon$,
  \[
    \phi_S(x)
    =\phi_S(1+\varepsilon)
     +S\,(x-1-\varepsilon)
     +\tfrac{\eta}{2}\bigl(x^2-(1+\varepsilon)^2\bigr).
  \]
\item[(d)] $\displaystyle\int_{I_S^+}\phi''_S(t)\,e^{\phi_S(t)}\,dt
  =S\bigl(1+O(S^{-2})\bigr)$.
\end{enumerate}
\end{lemma}

\begin{proof}
(a) is immediate from $\phi_S(x)=\eta x^2/2$ and $|x|\le 1$.

(b) For $x\in I_S^+$,
\[
  \phi_S(x)=\phi_S(1-\varepsilon)+\int_{1-\varepsilon}^x\phi'_S(t)\,dt
  \le \tfrac{\eta}{2}+2\varepsilon\,\sup_{I_S^+}\phi'_S
  \le \tfrac{\eta}{2}+2\varepsilon(\eta+S).
\]
Substituting $\eta=S^{-4}$, $\varepsilon=S^{-3}$ gives the bound
$O(S\varepsilon)=O(S^{-2})$. The bound $|e^{\phi_S}-1|\le \phi_S e^{\phi_S}
=O(S^{-2})$ follows. The bound on $I_S^-$ is identical by even symmetry.

(c) For $x>1+\varepsilon$, $\phi_S(x)-\phi_S(1+\varepsilon)
=\int_{1+\varepsilon}^x\phi'_S(t)\,dt=\int_{1+\varepsilon}^x(\eta t+S)\,dt$,
which evaluates to the stated expression.

(d) Using $e^{\phi_S}=1+O(S^{-2})$ uniformly on $I_S^+$ from (b) and
$\int_{I_S^+}\phi''_S\,dx=S+2\eta\varepsilon=S+O(S^{-7})$, we get
\[
  \int_{I_S^+}\phi''_S e^{\phi_S}\,dx
  =\int_{I_S^+}\phi''_S\,dx+\int_{I_S^+}\phi''_S\bigl(e^{\phi_S}-1\bigr)\,dx
  = S+O(S^{-7})+O(S^{-2})\cdot S
  = S\bigl(1+O(S^{-2})\bigr).\qedhere
\]
\end{proof}

\begin{lemma}\label{lem:normalization}
Under the choice \eqref{eq:specific-choice}, as $S\to\infty$:
\begin{enumerate}
\item[(a)] $\displaystyle\int_{1+\varepsilon}^{\infty}e^{-\phi_S(x)}\,dx
  =\frac{1}{S}+O(S^{-3})$.
\item[(b)] $\displaystyle Z_S = 2+\frac{2}{S}+O(S^{-3})$.
\item[(c)] Set
  $q_S:=\rho_S(E_S)=\rho_S\bigl(\{|x|>1+\varepsilon\}\bigr)$
  and $t_S:=\rho_S(T_S)$. Then
  \[
    q_S = \frac{1}{S}-\frac{1}{S^2}+O(S^{-3}),\qquad t_S=O(S^{-3}).
  \]
\end{enumerate}
\end{lemma}

\begin{proof}
(a) Substitute $u=x-1-\varepsilon$ in part (c) of
Lemma~\ref{lem:phi-bounds}, set $\widetilde S:=S+\eta(1+\varepsilon)=
S\bigl(1+O(\eta/S)\bigr)$, and use $\phi_S(1+\varepsilon)=O(S^{-2})$:
\[
  \int_{1+\varepsilon}^{\infty}e^{-\phi_S}\,dx
  = e^{-\phi_S(1+\varepsilon)}\int_0^{\infty}
     e^{-\widetilde S\,u-\eta u^2/2}\,du.
\]
The elementary bound $0\le 1-e^{-v}\le v$ for $v\ge0$ gives
\begin{equation}\label{eq:tail-expansion}
  0\le
  \frac{1}{\widetilde S}
  -\int_0^{\infty}e^{-\widetilde S\,u-\eta u^2/2}\,du
  \le
  \frac{\eta}{2}\int_0^\infty u^2e^{-\widetilde S u}\,du
  =\frac{\eta}{\widetilde S^3}.
\end{equation}
Since $\widetilde S=S+\eta(1+\varepsilon)$, this implies
\[
  \int_0^{\infty}e^{-\widetilde S\,u-\eta u^2/2}\,du
  =\frac{1}{S}+O\!\left(\frac{\eta}{S^2}+\frac{\eta}{S^3}\right)
  =\frac{1}{S}+O(S^{-6}).
\]
Combined with $e^{-\phi_S(1+\varepsilon)}=1+O(S^{-2})$ this gives
\[
  \int_{1+\varepsilon}^{\infty}e^{-\phi_S}\,dx
  =\bigl(1+O(S^{-2})\bigr)\bigl(1/S+O(S^{-6})\bigr)
  =\frac{1}{S}+O(S^{-3}),
\]
proving (a).

(b) Decompose $\int_0^{\infty}e^{-\phi_S}\,dx$ over the three regions
$[0,1-\varepsilon]$, $[1-\varepsilon,1+\varepsilon]$, $[1+\varepsilon,\infty)$.

On $[0,1-\varepsilon]$: $\phi_S(x)=\eta x^2/2$, so
$\int_0^{1-\varepsilon}e^{-\eta x^2/2}\,dx=(1-\varepsilon)+O(\eta)
=1-\varepsilon+O(S^{-4})$.

On $[1-\varepsilon,1+\varepsilon]$: $e^{-\phi_S}=1+O(S^{-2})$ by
Lemma~\ref{lem:phi-bounds}(b), so
$\int=2\varepsilon+O(\varepsilon/S^2)=2\varepsilon+O(S^{-5})$.

On $[1+\varepsilon,\infty)$: by part (a), $\int=1/S+O(S^{-3})$.

Adding and using $\varepsilon=S^{-3}$:
\[
  \int_0^{\infty}e^{-\phi_S}\,dx
  =1-\varepsilon+2\varepsilon+\frac{1}{S}+O(S^{-3})
  =1+\frac{1}{S}+O(S^{-3}).
\]
Multiplying by~$2$ (even symmetry) yields $Z_S=2+2/S+O(S^{-3})$.

(c) The layer mass is bounded by the Lebesgue measure of $T_S$ times
$\sup_{T_S}e^{-\phi_S}/Z_S\le 4\varepsilon/Z_S=O(\varepsilon)=O(S^{-3})$. For
$q_S$:
\[
  q_S=\frac{2}{Z_S}\int_{1+\varepsilon}^{\infty}e^{-\phi_S}\,dx
  =\frac{2/S+O(S^{-3})}{2+2/S+O(S^{-3})}
  =\frac{1}{S}\cdot\frac{1+O(S^{-2})}{1+1/S+O(S^{-3})}.
\]
Expanding the geometric series for $1/(1+1/S+O(S^{-3}))=
1-1/S+1/S^2+O(S^{-3})$:
\[
  q_S=\frac{1}{S}\Bigl(1-\frac{1}{S}+\frac{1}{S^2}+O(S^{-3})\Bigr)\bigl(1+O(S^{-2})\bigr)
   =\frac{1}{S}-\frac{1}{S^2}+O(S^{-3}).\qedhere
\]
\end{proof}

\subsection*{The test function $g_S$}

Define $g_S:\mathbb R\to[0,1]$ by
\begin{equation}\label{eq:gS}
  g_S(x):=
  \begin{cases}
    0 & x\in C_S,\\[4pt]
    \displaystyle
      \frac{\int_{1-\varepsilon}^{|x|}\phi''_S(t)e^{\phi_S(t)}\,dt}
           {\int_{I_S^+}\phi''_S(t)e^{\phi_S(t)}\,dt}
      & |x|\in (1-\varepsilon,1+\varepsilon),\\[12pt]
    1 & x\in E_S.
  \end{cases}
\end{equation}
On each of $C_S$ and $E_S$, $g_S$ is constant; on each layer it is (a
reparametrization of) the energy-minimizing interpolation between the
prescribed boundary values, as we verify in Lemma~\ref{lem:energy} below.

Continuity at the four endpoints $\pm 1\pm\varepsilon$ is built in: at
$|x|=1-\varepsilon$ the layer formula gives $g_S=0$, matching $C_S$; at
$|x|=1+\varepsilon$ it gives $g_S=1$, matching $E_S$. By construction $g_S$
is even, since the formula depends only on $|x|$.

Since $\phi''_S$ is bounded on $\mathbb R$ (it has supremum $\eta+(S/\varepsilon)
\|\kappa\|_\infty=O(S^4)$), $g_S'$ is bounded on each layer, hence $g_S$ is
globally Lipschitz. In particular, $g_S\in L^\infty(\mathbb R)\subset
L^2(\rho_S)$ and is locally Lipschitz.

\begin{lemma}\label{lem:energy}
The Brascamp--Lieb energy of $g_S$ is
\[
  \mathcal E_{\phi_S}(g_S)
  =\frac{2}{Z_S\int_{I_S^+}\phi''_S(t)e^{\phi_S(t)}\,dt}
  =\frac{2}{Z_S\,S}\bigl(1+O(S^{-2})\bigr).
\]
\end{lemma}

\begin{proof}
We first compute the energy of the optimal interpolant on a single layer.
Let $I=(a,b)$ be an interval and consider the variational problem
\begin{equation}\label{eq:single-J}
  J^*:=\inf\left\{\int_I\frac{(h'(x))^2}{\phi''_S(x)}\,e^{-\phi_S(x)}\,dx
        :\ h\in W^{1,2}(I),\ h(a)=0,\ h(b)=\Delta\right\},
\end{equation}
where the boundary values are interpreted as traces. Setting
$w(x):=e^{-\phi_S(x)}/\phi''_S(x)>0$, the integrand is $(h')^2 w$.
The Euler--Lagrange equation $(h'\,w)'=0$ gives $h'(x)=c/w(x)
=c\,\phi''_S(x)\,e^{\phi_S(x)}$ for some constant $c$. The boundary condition
$\int_a^b h'\,dx=\Delta$ forces
\[
  c=\frac{\Delta}{\int_I\phi''_S\,e^{\phi_S}\,dx}.
\]
Substituting $h'=c\phi''_S e^{\phi_S}$ into the integral gives
\[
  \int_I (h')^2 w\,dx
  =c^2\int_I\frac{(\phi''_S)^2 e^{2\phi_S}}{\phi''_S}\,e^{-\phi_S}\,dx
  =c^2\int_I\phi''_S e^{\phi_S}\,dx
  =\frac{\Delta^2}{\int_I\phi''_S e^{\phi_S}\,dx}.
\]
This is in fact the minimum: the integrand is a coercive non-negative
quadratic form in $h'$, with a strictly positive weight $w$ bounded above
and below on the compact interval $\overline I=[a,b]$ (here $w$ is bounded
above by $1/\eta$ and below by $\varepsilon/(S\|\kappa\|_\infty\,e^{O(S^{-2})})$,
both strictly positive). The set of admissible $h$ is the affine subspace
$\{h_0+W_0^{1,2}(I)\}$ with $h_0$ a fixed admissible interpolant, and $J$ is
strictly convex coercive on this subspace, so a minimizer exists and is
unique. The Euler--Lagrange formula above gives that minimizer.

By construction \eqref{eq:gS}, $g_S$ restricted to $I_S^+$ has $g_S(1-\varepsilon)
=0$, $g_S(1+\varepsilon)=1$, and
$g_S'(x)=\phi''_S(x)e^{\phi_S(x)}\big/\int_{I_S^+}\phi''_S e^{\phi_S}$,
which is exactly the Euler--Lagrange minimizer with $\Delta=1$. The same holds
on $I_S^-$ with $\Delta=-1$ (since $g_S$ ranges from $1$ at $-1-\varepsilon$
to $0$ at $-1+\varepsilon$, sign of $\Delta$ is opposite, but $\Delta^2=1$ in
either case). Outside the layers $g_S$ is constant, so $g_S'\equiv 0$ on
$C_S\cup E_S$. Therefore
\[
  \mathcal E_{\phi_S}(g_S)
  =\frac{1}{Z_S}\Biggl(\int_{I_S^+}\frac{(g_S')^2}{\phi''_S}e^{-\phi_S}\,dx
                 +\int_{I_S^-}\frac{(g_S')^2}{\phi''_S}e^{-\phi_S}\,dx\Biggr)
  =\frac{1}{Z_S}\bigl(J^*_{I_S^+}+J^*_{I_S^-}\bigr),
\]
and by the formula above with $\Delta^2=1$ on each layer and the symmetry
$\int_{I_S^+}\phi''_S e^{\phi_S}=\int_{I_S^-}\phi''_S e^{\phi_S}$,
\[
  \mathcal E_{\phi_S}(g_S)
  =\frac{2}{Z_S\int_{I_S^+}\phi''_S\,e^{\phi_S}\,dx}.
\]
The asymptotic statement now follows from
Lemma~\ref{lem:phi-bounds}(d).
\end{proof}

\begin{lemma}\label{lem:variance}
The variance of $g_S$ satisfies
\[
  \Var_{\rho_S}(g_S)
  = q_S(1-q_S)+O(S^{-3})
  = \frac{1}{S}-\frac{2}{S^2}+O(S^{-3}).
\]
\end{lemma}

\begin{proof}
On $E_S$, $g_S\equiv 1$; on $C_S$, $g_S\equiv 0$; on $T_S$, $0\le g_S\le 1$.
Hence
\[
  \int g_S\,d\rho_S
  = \rho_S(E_S)+\int_{T_S}g_S\,d\rho_S
  = q_S+\mathcal R_1,\qquad |\mathcal R_1|\le t_S=O(S^{-3}),
\]
\[
  \int g_S^2\,d\rho_S
  = \rho_S(E_S)+\int_{T_S}g_S^2\,d\rho_S
  = q_S+\mathcal R_2,\qquad |\mathcal R_2|\le t_S=O(S^{-3}),
\]
by Lemma~\ref{lem:normalization}(c). Therefore
\[
  \Var_{\rho_S}(g_S)
  =(q_S+\mathcal R_2)-(q_S+\mathcal R_1)^2
  = q_S-q_S^2+\bigl(\mathcal R_2-2q_S\mathcal R_1-\mathcal R_1^2\bigr)
  = q_S(1-q_S)+O(S^{-3}),
\]
where in the last step we used $q_S=O(S^{-1})$ to absorb $-2q_S\mathcal R_1
=O(S^{-4})$ and $\mathcal R_1^2=O(S^{-6})$ into the $O(S^{-3})$ error.

Substituting $q_S=1/S-1/S^2+O(S^{-3})$ from
Lemma~\ref{lem:normalization}(c),
\[
  q_S^2
  =\frac{1}{S^2}\Bigl(1-\frac{1}{S}+O(S^{-2})\Bigr)^{\!2}
  =\frac{1}{S^2}-\frac{2}{S^3}+O(S^{-4}),
\]
hence
\[
  q_S(1-q_S)=q_S-q_S^2
  =\Bigl(\tfrac{1}{S}-\tfrac{1}{S^2}\Bigr)-\Bigl(\tfrac{1}{S^2}-\tfrac{2}{S^3}\Bigr)+O(S^{-3})
  =\tfrac{1}{S}-\tfrac{2}{S^2}+O(S^{-3}). \qedhere
\]
\end{proof}

\section{\texorpdfstring{The spectral gap above
$\mathcal O_{\mathrm{BL},\phi_S}$ closes}{The spectral gap above OBL closes}}

Let
\[
  c_S := \int_{\mathbb R} g_S\,d\rho_S,\qquad
  f_S := g_S - c_S.
\]
Since $g_S$ is even and $\rho_S$ is even, $c_S\in\mathbb R$ and $f_S$ is even
with $\rho_S$-mean zero. The Brascamp--Lieb deficit and the variance are
unchanged by additive constants, so
\begin{equation}\label{eq:fS-vs-gS}
  \mathcal E_{\phi_S}(f_S)=\mathcal E_{\phi_S}(g_S),\qquad
  \Var_{\rho_S}(f_S)=\Var_{\rho_S}(g_S),\qquad
  \delta_{\mathrm{BL},\phi_S}(f_S)=\delta_{\mathrm{BL},\phi_S}(g_S).
\end{equation}

\begin{lemma}\label{lem:orth}
$f_S$ is orthogonal in $L^2(\rho_S)$ to the optimizer space
$\mathcal O_{\mathrm{BL},\phi_S}=\mathrm{span}_{\mathbb R}\{1,\phi'_S\}$.
Consequently
\[
  \dist_{L^2(\rho_S)}(f_S,\mathcal O_{\mathrm{BL},\phi_S})^2
  =\|f_S\|_{L^2(\rho_S)}^2
  =\Var_{\rho_S}(g_S).
\]
\end{lemma}

\begin{proof}
$\langle f_S,1\rangle_{L^2(\rho_S)}=\int f_S\,d\rho_S=0$ by construction. For
the second generator, since $f_S$ is even, $\phi'_S$ is odd, and $e^{-\phi_S}$
is even, the integrand $f_S(x)\phi'_S(x)e^{-\phi_S(x)}$ is an odd function in
$L^1(\mathbb R)$, so
\[
  \langle f_S,\phi'_S\rangle_{L^2(\rho_S)}
  =\frac{1}{Z_S}\int_{\mathbb R}f_S(x)\phi'_S(x)e^{-\phi_S(x)}\,dx
  =0.
\]
Since $\phi'_S$ is non-constant (it ranges over an interval containing
$\pm S$), the two generators $1$ and $\phi'_S$ are linearly independent in
$L^2(\rho_S)$, so $f_S\perp\mathcal O_{\mathrm{BL},\phi_S}$. The distance
formula follows from $\|f_S\|^2=\int f_S^2\,d\rho_S=\Var_{\rho_S}(g_S)$,
since $f_S$ has $\rho_S$-mean zero.
\end{proof}

\begin{proposition}\label{prop:asymptotics}
With the choice \eqref{eq:specific-choice}, as $S\to\infty$,
\[
  \mathcal E_{\phi_S}(f_S) = \frac{1}{S}-\frac{1}{S^2}+O(S^{-3}),
  \qquad
  \Var_{\rho_S}(f_S) = \frac{1}{S}-\frac{2}{S^2}+O(S^{-3}),
\]
and consequently
\[
  \delta_{\mathrm{BL},\phi_S}(f_S)
  = \frac{1}{S^2}+O(S^{-3}).
\]
In particular $\delta_{\mathrm{BL},\phi_S}(f_S)>0$ for all sufficiently
large~$S$, and the Brascamp--Lieb Rayleigh quotient
$\mathcal E_{\phi_S}(f_S)/\Var_{\rho_S}(f_S)$ converges to~$1$ from above at
rate $1/S$.
\end{proposition}

\begin{proof}
Combining Lemma~\ref{lem:energy} with the identity
$2/(Z_S\,S)=q_S\cdot\bigl(1+O(S^{-2})\bigr)$ from
Lemma~\ref{lem:normalization}(c) (alternatively, computing directly:
$2/(Z_S S)=1/(S+1+O(S^{-2}))=1/S-1/S^2+O(S^{-3})$),
\[
  \mathcal E_{\phi_S}(g_S)
  =\frac{2}{Z_S\,S}\bigl(1+O(S^{-2})\bigr)
  =\Bigl(\tfrac{1}{S}-\tfrac{1}{S^2}+O(S^{-3})\Bigr)\bigl(1+O(S^{-2})\bigr)
  =\tfrac{1}{S}-\tfrac{1}{S^2}+O(S^{-3}).
\]
Lemma~\ref{lem:variance} gives the variance asymptotics. Subtracting,
\[
  \delta_{\mathrm{BL},\phi_S}(f_S)
  =\mathcal E_{\phi_S}(g_S)-\Var_{\rho_S}(g_S)
  =\Bigl(\tfrac{1}{S}-\tfrac{1}{S^2}\Bigr)
   -\Bigl(\tfrac{1}{S}-\tfrac{2}{S^2}\Bigr)+O(S^{-3})
  =\tfrac{1}{S^2}+O(S^{-3}).
\]
The displayed asymptotic implies $\delta_{\mathrm{BL},\phi_S}(f_S)>0$ for all
sufficiently large~$S$.
\end{proof}

\begin{theorem}[1D part of Theorem~\ref{thm:l2-fails}]\label{thm:1d}
With the choice \eqref{eq:specific-choice} and the test function $f_S$
defined above,
\[
  \frac{\dist_{L^2(\rho_S)}(f_S,\mathcal O_{\mathrm{BL},\phi_S})^2}
       {\delta_{\mathrm{BL},\phi_S}(f_S)}
  =S\bigl(1+O(S^{-1})\bigr)\xrightarrow[S\to\infty]{}+\infty.
\]
Therefore the constant in any putative inequality
\eqref{eq:would-be-stability} must blow up at rate at least~$\sqrt S$.
\end{theorem}

\begin{proof}
By Lemma~\ref{lem:orth} and Lemma~\ref{lem:variance},
\[
  \dist_{L^2(\rho_S)}(f_S,\mathcal O_{\mathrm{BL},\phi_S})^2
  =\Var_{\rho_S}(g_S)
  =\tfrac{1}{S}-\tfrac{2}{S^2}+O(S^{-3}).
\]
Combined with $\delta_{\mathrm{BL},\phi_S}(f_S)=1/S^2+O(S^{-3})$ from
Proposition~\ref{prop:asymptotics},
\[
  \frac{\dist^2}{\delta_{\mathrm{BL}}}
  =\frac{1/S-2/S^2+O(S^{-3})}{1/S^2+O(S^{-3})}
  =\frac{(1/S)\bigl(1-2/S+O(S^{-2})\bigr)}
        {(1/S^2)\bigl(1+O(S^{-1})\bigr)}
  =S\bigl(1+O(S^{-1})\bigr).
\]
Equivalently $\dist/\delta_{\mathrm{BL}}^{1/2}=\sqrt S\,\bigl(1+O(S^{-1})\bigr)$,
so any constant in \eqref{eq:would-be-stability} must satisfy
$C\ge\sqrt S\,(1+o(1))$.
\end{proof}

\subsection*{Higher dimensions}

Fix $d\ge 2$ and define the product potential
\[
  \Phi_S(x,y):=\phi_S(x)+\frac{|y|^2}{2},
  \qquad (x,y)\in\mathbb R\times\mathbb R^{d-1}.
\]
Then $\Phi_S$ is $C^\infty$, strictly convex (as a sum of strictly convex
functions in disjoint variables), with
$D^2\Phi_S=\mathrm{diag}(\phi''_S(x),I_{d-1})$ and
$(D^2\Phi_S)^{-1}=\mathrm{diag}(1/\phi''_S(x),I_{d-1})$. The associated Gibbs
measure is the product
$\rho_{\Phi_S}=\rho_S\otimes\gamma_{d-1}$, with $\gamma_{d-1}$ the standard
Gaussian on $\mathbb R^{d-1}$. The optimizer space is
\[
  \mathcal O_{\mathrm{BL},\Phi_S}
  =\bigl\{a\cdot\nabla\Phi_S+b: a\in\mathbb R^d,b\in\mathbb R\bigr\}
  =\mathrm{span}_{\mathbb R}\bigl\{1,\,\phi'_S(x),\,y_1,\dots,y_{d-1}\bigr\},
\]
a $(d+1)$-dimensional subspace of $L^2(\rho_{\Phi_S})$.

Set $F_S(x,y):=f_S(x)$, so $\nabla F_S=(f_S'(x),0_{\mathbb R^{d-1}})$ and
$(D^2\Phi_S)^{-1}\nabla F_S=(f_S'(x)/\phi''_S(x),0)$. Therefore
\[
  \bigl\langle (D^2\Phi_S)^{-1}\nabla F_S,\nabla F_S\bigr\rangle
  =\frac{(f_S'(x))^2}{\phi''_S(x)}
  \qquad\text{pointwise}.
\]
Fubini and $\int d\gamma_{d-1}=1$ give
\[
  \int_{\mathbb R^d}\bigl\langle (D^2\Phi_S)^{-1}\nabla F_S,\nabla F_S\bigr\rangle
   \,d\rho_{\Phi_S}
  =\int_{\mathbb R}\frac{(f_S')^2}{\phi''_S}\,d\rho_S
  =\mathcal E_{\phi_S}(f_S),
\]
\[
  \int F_S\,d\rho_{\Phi_S}=\int f_S\,d\rho_S=0,\qquad
  \Var_{\rho_{\Phi_S}}(F_S)=\int f_S^2\,d\rho_S=\Var_{\rho_S}(f_S),
\]
hence
$\delta_{\mathrm{BL},\Phi_S}(F_S)=\delta_{\mathrm{BL},\phi_S}(f_S)$.

Orthogonality of $F_S$ to each generator of $\mathcal O_{\mathrm{BL},\Phi_S}$:
\[
  \langle F_S,1\rangle_{L^2(\rho_{\Phi_S})}=\int f_S\,d\rho_S=0,
\]
\[
  \langle F_S,\phi'_S\rangle_{L^2(\rho_{\Phi_S})}
  =\Bigl(\int f_S\,\phi'_S\,d\rho_S\Bigr)\Bigl(\int d\gamma_{d-1}\Bigr)
  =0\cdot 1=0
  \quad\text{by Lemma~\ref{lem:orth}},
\]
\[
  \langle F_S,y_j\rangle_{L^2(\rho_{\Phi_S})}
  =\Bigl(\int f_S\,d\rho_S\Bigr)\Bigl(\int y_j\,d\gamma_{d-1}\Bigr)
  =0\cdot 0=0,
  \quad j=1,\dots,d-1.
\]
Hence $F_S\perp\mathcal O_{\mathrm{BL},\Phi_S}$ in $L^2(\rho_{\Phi_S})$, so
\[
  \dist_{L^2(\rho_{\Phi_S})}(F_S,\mathcal O_{\mathrm{BL},\Phi_S})^2
  =\|F_S\|_{L^2(\rho_{\Phi_S})}^2
  =\|f_S\|_{L^2(\rho_S)}^2,
\]
and the divergent ratio of Theorem~\ref{thm:1d} is unchanged when one passes
from $(\phi_S,f_S)$ to $(\Phi_S,F_S)$. This proves
Theorem~\ref{thm:l2-fails} in every dimension $d\ge 1$.\hfill$\square$

\begin{remark}[On the failure of the Poincar\'e route]\label{rem:affine}
The Brascamp--Lieb Dirichlet form is invariant under affine reparametrizations
of the underlying space: for $T(x)=Ax+b$ with $A\in GL_d(\mathbb R)$, set
$\widetilde\phi:=\phi\circ T^{-1}$ and $\widetilde f:=f\circ T^{-1}$. Then
$\rho_{\widetilde\phi}=T_\#\rho_\phi$, $D^2\widetilde\phi(y)=
A^{-T}D^2\phi(T^{-1}y)A^{-1}$, $\nabla\widetilde f(y)=A^{-T}\nabla f(T^{-1}y)$,
and a direct computation yields
\[
  \mathcal E_{\widetilde\phi}(\widetilde f)=\mathcal E_\phi(f),\qquad
  \Var_{\rho_{\widetilde\phi}}(\widetilde f)=\Var_{\rho_\phi}(f),\qquad
  \mathcal O_{\mathrm{BL},\widetilde\phi}=\mathcal O_{\mathrm{BL},\phi}\circ T^{-1}.
\]
The pushforward $L^2$-distance to the optimizer space is also preserved:
$\dist_{L^2(\rho_{\widetilde\phi})}(\widetilde f,\mathcal O_{\mathrm{BL},
\widetilde\phi})=\dist_{L^2(\rho_\phi)}(f,\mathcal O_{\mathrm{BL},\phi})$.
Therefore the ratio
$\dist_{L^2}^2/\delta_{\mathrm{BL}}$ is an affine invariant of $(\phi,f)$.

The Euclidean Poincar\'e constant $\lambda_1(\rho_\phi)^{-1}$, by contrast,
\emph{is not} affinely invariant: stretching a measure by an anisotropic
diagonal matrix multiplies it by the squared anisotropy ratio. Any attempt to
make the would-be $L^2$-stability constant blow up by exhibiting log-concave
measures with unbounded ordinary Poincar\'e constant (such as long thin
rectangles) is therefore doomed: the stretching that creates the large
Poincar\'e constant is precisely the affine map under which the
Brascamp--Lieb deficit, the $L^2$-distance to the optimizer space, and the
optimizer space itself transform covariantly, leaving the relevant ratio
unchanged.

The construction above evades this obstruction by creating a genuinely new
near-eigenmode of the Brascamp--Lieb operator $L_{\phi_S}$, in the orthogonal
complement of $\mathcal O_{\mathrm{BL},\phi_S}$ rather than within an affine
copy of it. Concretely, the even function $f_S$ realizes a Rayleigh quotient
$\mathcal E_{\phi_S}(f_S)/\Var_{\rho_S}(f_S)\to 1$ even though $f_S$ stays
orthogonal to the entire optimizer space $\{1,\phi'_S\}$; this reflects the
spectral collapse described after Theorem~\ref{thm:l2-fails} and is invariant
under affine reparametrization.
\end{remark}

\begin{thebibliography}{10}
\bibitem{BLpaper}
J.~M. Machado and J.~P. G. Ramos,
\textit{Quantitative stability for Brascamp--Lieb and moment measures},
arXiv:2511.22636v2.

\bibitem{BL76}
H.~J. Brascamp and E.~H. Lieb,
\textit{On extensions of the Brunn--Minkowski and Pr\'ekopa--Leindler theorems,
including inequalities for log-concave functions, and with an application to
the diffusion equation},
J.~Funct.~Anal. 22 (1976), 366--389.
\end{thebibliography}

\end{document}
